\begin{problem}[理论考试][绝热容器与热力学过程]{27}
  如图所示，某绝热容器被二块装有阀门 $\mathrm{K}_{1}$ 和 $\mathrm{K}_{2}$ 的固定绝热隔板分隔成相等体积 $V_{0}$ 的三室A、B、C，$V_{\mathrm{A}} = V_{\mathrm{B}} = V_{\mathrm{C}} = V_{0}$。容器左端用绝热活塞H封闭，左侧A室装有 $\nu_{1} = 1$ 摩尔单原子分子气体，处在压强为 $P_{0}$、温度为 $T_{0}$ 的平衡态；中段B室为真空；右侧C室装有 $\nu_{2} = 2$ 摩尔双原子分子气体，测得其平衡态温度为 $T_{c} = 0.50T_{0}$。初始时刻 $\mathrm{K}_{1}$ 和 $\mathrm{K}_{2}$ 都处在关闭状态。然后系统依次经历如下与外界无热量交换的热力学过程：
  \begin{subproblem}
    打开 $K_{1}$，让 $V_{A}$ 中的气体自由膨胀到中段真空 $V_{B}$ 中；等待气体达到平衡态时，缓慢推动活塞 H 压缩气体，使得 A 室体积减小了 30\%（$V_{A}' = 0.70V_{0}$）。求压缩过程前后，该部分气体的平衡态温度及压强；
  \end{subproblem}
  \begin{subproblem}
    保持 $\mathrm{K}_{1}$ 开启，打开 $\mathrm{K}_{2}$，让容器中的两种气体自由混合后共同达到平衡态。求此时混合气体的温度和压强；
  \end{subproblem}
  \begin{subproblem}
    保持 $\mathrm{K}_{1}$ 和 $\mathrm{K}_{2}$ 同时处在开启状态，缓慢拉动活塞 $\mathrm{H}$，使得 $\mathrm{A}$ 室体积恢复到初始体积 $V_{\mathrm{A}}^{\prime \prime} = V_{0}$。求此时混合气体的温度和压强。
  \end{subproblem}
  提示：上述所有过程中，气体均可视为理想气体，计算结果可含数值的指数式或分式。根据热力学第二定律，当一种理想气体构成的热力学系统从初态 $(p_i, T_i, V_i)$ 经过一个绝热可逆过程（准静态绝热过程）到达终态 $(p_f, T_f, V_f)$ 时，其状态参数满足方程：
  \begin{equation}
    \left(\Delta S_{1}\right)_{i f} = \nu_{1} C_{V_{1}} \ln (\frac{T_{f}}{T_{i}}) + \nu_{1} R \ln (\frac{V_{f}}{V_{i}}) = 0
    \label{eq:I}
  \end{equation}
  其中，$\nu_{1}$ 为该气体的摩尔数，$C_{V_{1}}$ 为它的定容摩尔热容量，R 为普适气体常量；而当热力学系统由二种理想气体组成，则方程（I）需修改为
  \begin{equation}
    \left(\Delta S_{1}\right)_{i f} + \left(\Delta S_{2}\right)_{i f} = 0
    \label{eq:II}
  \end{equation}
\end{problem}